Regression estimators of the global average treatment effect under network interference fit the outcome on a unit's own treatment and its neighbours' treatments, then predict with everyone treated and everyone in control. Their error is understood when the fitted model is correct; we ask what determines it otherwise. Because randomised exposure never reaches those two corners, regression estimation is extrapolation to the corner: its error is a variance that grows as the square root of the degree plus a bias, set by the curvature of the true response, that does not shrink with the sample size. On synthetic graphs and two real Facebook graphs with simulated outcomes, the variance equals the corner leverage prediction and the bias an exact projection bias computed from the graph and design alone (correlation 0.9996 over 320 cells). Regression beats exact Horvitz--Thompson at every degree from two upward, and loses to relaxed-exposure estimators or to ignoring interference only for step-like responses. Data-only scores of the bias do not beat the corner distance, except one that assumes the response does not depend on degree and fails when it does, and a bias-aware interval attains coverage only by being wide.
